% Results tables (I, III--V); calculation inventory (II) is in Methods.

\begin{table}[!ht]
\caption{\label{tab:I}Archived tetramer cluster map (legacy PBE without D3). Cartesian coordinates are \emph{not} affinely scaled across vacuum-box labels; these are not monolayer strain coefficients and are not mixed with periodic PBE+D3 $\mathcal{S}$. The valid ionic-relaxation test is periodic $n{=}1$ $P$ (Fig.~\ref{fig:synergy}(b)), not Table~\ref{tab:relax}.
$E_{\mathrm{sub}}=(E_{\mathrm{doped}}-E_{\mathrm{pristine}})/n_{\mathrm{dopants}}$ at $\epsilon{=}0$ with no chemical-potential reference (not a standard formation enthalpy).
$n_{\mathrm{dopants}}$ is counted from archived input coordinates (nominal 5\% target; B: 13, N: 12, P: 12).}
\begin{ruledtabular}
\begin{tabular}{lccccc}
System & $n_{\mathrm{dop}}$ & $E(\epsilon{=}0)$ ($E_{\mathrm{h}}$) & $E(\epsilon{=}+5\%)$ ($E_{\mathrm{h}}$) & $\alpha$ (meV/\%) & $E_{\mathrm{sub}}$ (eV/dopant) \\
\hline
Pristine & --- & $-1363.5512$ & $-1363.5509$ & $-1.0$ & --- \\
B-doped & 13 & $-1328.4900$ & $-1328.4090$ & $+56.9$ & $+73.4$ \\
N-doped & 12 & $-1414.3812$ & $-1414.4053$ & $-303.2$ & $-115.3$ \\
P-doped & 12 & $-1370.4965$ & $-1370.4888$ & $-19.1$ & $-15.8$ \\
\end{tabular}
\end{ruledtabular}
\end{table}

\begin{table}[!ht]
\caption{\label{tab:relaxext}Cluster-map diagnostic for the tetramer $P$-substitution four-corner grid (pristine and $P$ at two vacuum-box labels). Cartesian coordinates are \emph{not} affinely scaled; only the vacuum box changes, so this is \emph{not} a monolayer strain-retention test. The valid ionic-relaxation control is periodic $n{=}1$ $P$ at affine $+3$\%: $\mathcal{S}_{\mathrm{rigid}}=-31.9$~meV/atom $\to$ $\mathcal{S}_{\mathrm{relaxed}}\approx 0$. Relaxed energies: PBE+D3 fixed-cell geometry optimization (400~Ry) from archived inputs in the open repository. $\mathcal{S}$ follows Eq.~\eqref{eq:synergy_order} with $N{=}240$ atoms.}
\begin{ruledtabular}
\begin{tabular}{lcc}
Quantity & Rigid (legacy tetramer) & Relaxed (geometry opt.) \\
\hline
$E$(pristine, $\epsilon{=}0$) ($E_{\mathrm{h}}$) & $-1363.5512$ & $-1363.9551$ \\
$E$(pristine, $\epsilon{=}3$\%) ($E_{\mathrm{h}}$) & $-1363.5522$ & $-1363.9550$ \\
$E$(P, $\epsilon{=}0$) ($E_{\mathrm{h}}$) & $-1370.4965$ & $-1374.4415$ \\
$E$(P, $\epsilon{=}3$\%) ($E_{\mathrm{h}}$) & $-1370.4890$ & $-1374.4615$ \\
$\mathcal{S}$ (meV/atom) & $+0.96$ & $-2.28$ \\
Sign of $\mathcal{S}$ vs.\ rigid & $+$ & $-$ (not preserved) \\
\end{tabular}
\end{ruledtabular}
\end{table}

\begin{table}[!ht]
\caption{\label{tab:sgrid}Cluster-map diagnostic on two independently rerolled tetramer substitution maps (\emph{not} a monolayer strain-placement scan). Cartesian coordinates are identical across archived vacuum-box labels; only the vacuum box and dopant identities change.
The \emph{reference} cluster map supplies the archived $\alpha$ in Table~\ref{tab:I}; those slopes and $\mathcal{S}_{+3\%}^{\mathrm{tet}}$ are cluster diagnostics on the legacy-PBE grid (no D3), not monolayer strain coefficients.
The \emph{alternate} map repeats the same vacuum-box labels ($\epsilon\in\{-5,-2.5,0,+2.5,+3,+5\}$\%; 18 rigid single-points) with PBE+D3 (Sec.~\ref{sec:methods_s4_seed}).
$\mathcal{S}_{+3\%}^{\mathrm{tet}}$ follows Eq.~\eqref{eq:synergy_order} on the archived tetramer four-corner grid ($N{=}240$).
Periodic $\mathcal{S}(n)$ is tabulated in the main text; $n{=}4$ ranking $P{>}B{>}N$ is reproduced on periodic seeds 42, 137, and 271.}
\begin{ruledtabular}
\begin{tabular}{lcccc}
Dopant & Placement & $\alpha$ (meV/\%) & $\mathcal{S}_{+3\%}^{\mathrm{tet}}$ (meV/atom) & Status \\
\hline
B & reference & $+56.9$ & $+3.8$ & cluster diagnostic (Table~\ref{tab:placement}) \\
N & reference & $-303.2$ & $+5.7$ & cluster diagnostic (Table~\ref{tab:placement}) \\
P & reference & $-19.1$ & $+1.0$ & cluster diagnostic (Table~\ref{tab:placement}) \\
B & alternate (seed~137) & $-45$ & $+1.9$ & cluster diagnostic (Table~\ref{tab:placement}) \\
N & alternate (seed~137) & $-60$ & $+0.3$ & cluster diagnostic (Table~\ref{tab:placement}) \\
P & alternate (seed~137) & $+1028$ & $-32.3$ & cluster diagnostic (Table~\ref{tab:placement}) \\
\end{tabular}
\end{ruledtabular}
\end{table}

\begin{table}[!ht]
\caption{\label{tab:V}Local dopant--cage structure from archived tetramer XYZ (mean nearest C distance $\bar{d}$ in \AA; standard deviation $\sigma(\bar{d})$ across dopant sites; covalent-radius excess vs.\ C in pm).
The Cartesian coordinates are \emph{not} affinely scaled between the labeled $\epsilon{=}0$ and $+3$\% tags, so $\Delta\bar{d}=\bar{d}(\mathrm{tag}{=}+3\%)-\bar{d}(\mathrm{tag}{=}0)$ is a cluster-map / vacuum-box diagnostic, not tensile-load coupling (Discussion mechanistic synthesis).
Values are computed from archived tetramer coordinates in the open data release; Mayer bond order and Bader charges are not included.}
\begin{ruledtabular}
\begin{tabular}{lccccc}
Dopant & $\bar{d}(\epsilon{=}0)$ (\AA) & $\bar{d}(\epsilon{=}+3\%)$ (\AA) & $\Delta\bar{d}$ (\AA) & $\sigma(\bar{d})$ @ $+3$\% (\AA) & $\Delta r_{\mathrm{cov}}$ (pm) \\
\hline
B & 1.384 & 1.393 & $+0.0085$ & 0.019 & $+7$ \\
N & 1.393 & 1.402 & $+0.0087$ & 0.024 & $-6$ \\
P & 1.384 & 1.385 & $+0.00003$ & 0.0004 & $+30$ \\
\end{tabular}
\end{ruledtabular}
\end{table}
